% =============================================================================
%  Geometry or Dynamics? What Diffusion Maps Recover from a Stochastic
%  Double-Well System
%
%  Build:   pdflatex main
%           bibtex  main
%           pdflatex main
%           pdflatex main
%
%  Figures are read from ../outputs/figures/ (see \graphicspath below), so the
%  experiment must have been run at least once (python run_experiment.py).
% =============================================================================

\documentclass[11pt]{article}

\usepackage[letterpaper,margin=1in]{geometry}
\usepackage{amsmath,amssymb,amsthm}
\usepackage{graphicx}
\usepackage{booktabs}
\usepackage{caption}
\usepackage{subcaption}
\usepackage{algorithm}
\usepackage{algpseudocode}
\usepackage[round]{natbib}
\usepackage{xcolor}
\usepackage[colorlinks=true,linkcolor=blue!50!black,citecolor=blue!50!black,
            urlcolor=blue!50!black]{hyperref}
\usepackage{url}

\graphicspath{{../outputs/figures/}{figures/}}

\newtheorem{proposition}{Proposition}
\newtheorem{remark}{Remark}

\newcommand{\R}{\mathbb{R}}
\newcommand{\dd}{\,\mathrm{d}}
\newcommand{\eps}{\varepsilon}
\newcommand{\Kt}{\widetilde{K}}
\newcommand{\norm}[1]{\left\lVert #1 \right\rVert}
\newcommand{\diag}{\operatorname{diag}}
\newcommand{\sgn}{\operatorname{sign}}
\DeclareMathOperator{\Var}{Var}

\title{\bf Geometry or Dynamics? What Diffusion Maps Recover\\
from a Stochastic Double-Well System}

\author{A.\ L.\ Jackson}

\date{\today}

\begin{document}
\maketitle

% =============================================================================
\begin{abstract}
% =============================================================================
Diffusion maps are routinely used to extract slow reaction coordinates from
high-dimensional observations of stochastic systems. We ask, on a problem where
the ground truth is known exactly, what they actually recover. The hidden state
is the overdamped Langevin dynamics of a particle in a symmetric quartic double
well, $\dd X_t = (X_t - X_t^3)\dd t + \sigma\dd W_t$, integrated by
Euler--Maruyama over $T = 600$ time units --- about $49$ Kramers escape times,
yielding $61$ committed barrier crossings --- and subsampled in time to $N=3000$
points so that a dense $O(N^2)$ kernel remains affordable. Two observation models
lift the scalar state into $\R^{15}$ through a random Gaussian map with additive
noise: a \emph{polynomial} map whose image is a gently bent arc, and an
\emph{RBF sensor-array} map whose image is strongly curved.

Two findings separate cleanly. First, on the question of \emph{state recovery},
the diffusion map decisively beats PCA where the manifold is curved and ties
where it is not: on the RBF observations
$\lvert\rho_{\mathrm{Spearman}}\rvert = 0.9998$ against PCA's $0.6624$ (PC$_1$
is a non-monotone hump that confuses $X \approx -0.3$ with $X \approx 1.5$),
while on the polynomial observations both reach $0.9997$. The advantage is not
free: it survives only at small bandwidth, and as $\eps$ grows the diffusion map
degrades monotonically \emph{onto} the PCA baseline ($0.9998 \to 0.654$ over a
$128\times$ sweep), which explains why a median-distance bandwidth heuristic
makes the two methods look identical.

Second, on the question of \emph{metastability}, the normalization controls how
strongly the sampling density influences the answer, while the observation
geometry controls how much that matters. Sweeping $\alpha$ leaves state recovery
untouched ($\lvert\rho\rvert = 0.9998$ throughout) while polynomial-model well
accuracy falls from $0.9960$ at $\alpha=0$ to $0.5804$ --- poor, though still
above the $0.5$ chance baseline --- at $\alpha=1$. On the RBF model the wells are
already separated geometrically, so the effect is much weaker. The spectra also
separate density-sensitive and density-free regimes: at $\alpha=1$ the RBF gap
ratios sit near $(4,9)$, the $k^2$ law of the Laplace--Beltrami operator on an
interval, while at $\alpha=1/2$ the leading gap grows as the physical barrier is
deepened. Because the nonlinear observation map changes the metric, the latter
is a dynamics-informed proxy rather than an exact estimate of the SDE generator.

We also report a metric failure worth guarding against: at very small bandwidth
the kernel becomes nearly reducible, and the resulting eigenvector --- effectively
a two-point component indicator, constant elsewhere to within $4\times10^{-8}$ ---
still returns $\lvert\rho_{\mathrm{Spearman}}\rvert = 0.9997$. Rank correlation is
scale-free and cannot tell that the ordering lies eight orders of magnitude below
the coordinate's full range. We use a conservative MST-scale-constrained
bandwidth scan and diagnostics for near-unit eigenvalues and coordinate collapse.
The practical rule is to choose $\alpha$ from the question, inspect the spectrum,
and never read a diffusion-map score without checking kernel conditioning.
\end{abstract}

\tableofcontents
\newpage

% =============================================================================
\section{Introduction}
% =============================================================================

\subsection{Problem statement}

A recurring situation in the physical sciences is that a system's essential
behavior is governed by a small number of slow degrees of freedom, while the
data actually collected are numerous, indirect, and entangled. A protein folds
along what is effectively one or two reaction coordinates, but the experiment
returns thousands of atomic positions. A climate oscillation is
low-dimensional, but the observations are gridded fields. The scientific task is
inverse: given the high-dimensional measurements alone, and no knowledge of the
map that produced them, recover the underlying slow variable.

We study that task where the ground truth is known exactly, so recovery can be
measured rather than asserted. The hidden state is the position $X_t \in \R$ of
an overdamped Brownian particle in a symmetric quartic double-well potential ---
the canonical minimal model of a two-state system with rare transitions. The
observations are produced by a deliberately obstructive pipeline: a nonlinear
feature map, a random linear lift into $\R^{15}$, additive Gaussian noise, and
column standardization, so that no measured coordinate is recognizably $X$.

Two questions are asked, and it matters that they are distinct:

\begin{quote}
\textbf{Q1 (geometry).} Can an unsupervised manifold-learning method recover the
hidden scalar state from the high-dimensional encoding, and does it do so better
than a linear method?

\textbf{Q2 (dynamics).} Does the leading coordinate recover the \emph{slow}
variable --- the one tracking which well the system occupies --- so that the two
metastable states are separated?
\end{quote}

The pre-registered hypothesis was that both would be answered affirmatively by
the same object, the leading nontrivial diffusion coordinate, because diffusion
maps preserve local geometric connectivity. Diffusion maps
\citep{coifman2006diffusion} are the natural candidate: they build a random walk
on the data whose slow eigenmodes parameterize the underlying geometry, and there
is theory \citep{nadler2006diffusion,coifman2008diffusion} connecting those modes
to reaction coordinates of stochastic systems. \citet{banisch2017diffusion}
extend the construction to the generators of arbitrary non-degenerate It\^o
diffusions --- exactly this setting. PCA \citep{jolliffe2002principal} is the
natural linear baseline.

\subsection{Summary of findings}

The two questions receive different answers, and separating them is the main
content of this report.

\begin{enumerate}
\item \textbf{Q1: yes, and curvature strongly affects the margin.} On RBF
sensor-array observations the diffusion map
attains $\lvert\rho_{\mathrm{Spearman}}\rvert = 0.9998$ against PCA's $0.6624$;
on polynomial observations the two tie at $0.9997$. The difference is not the
method but the geometry: the polynomial feature map leaves a dominant monotone
imprint of $X$ (one linear direction carries $64\%$ of the variance), so a chord
orders the arc as well as its arclength does. A comparison run only on the
polynomial model --- as in our own preliminary configuration --- would have
concluded, wrongly, that diffusion maps offer nothing over PCA
(Sections~\ref{sec:res-q1}, \ref{sec:disc-pilot}).

\item \textbf{Q1, caveat: the advantage lives at small bandwidth.} Sweeping
$\eps$ over $128\times$, RBF recovery decays monotonically from $0.9998$ onto
PCA's $0.662$. The diffusion map degenerates to a linear method as
$\eps \to \infty$, so a bandwidth heuristic that overshoots destroys precisely
the nonlinearity it was meant to exploit (Section~\ref{sec:res-bandwidth}).

\item \textbf{Q2: controlled jointly by $\alpha$ and observation geometry.} The
observations are a memoryless function of $X$, so sampling density is the only
trace of the dynamics available to the point cloud. Sweeping $\alpha$ leaves
state recovery flat at $0.9998$ while polynomial-model well accuracy falls from
$0.9960$ ($\alpha=0$) to $0.5804$ ($\alpha=1$), whereas the geometrically
separated RBF model changes little (Section~\ref{sec:res-alpha}).

\item \textbf{The spectra identify which operator was computed.} At $\alpha=1$
the gap ratios $(1-\lambda_k)/(1-\lambda_1)$ measure $(4.19,\,8.38)$ against the
$(4,\,9)$ of the Laplace--Beltrami operator on an interval, and remain pinned
there across every barrier depth --- geometry is blind to the dynamics. At
$\alpha=1/2$ they instead rise from $10.2$ to $33.0$ as the barrier deepens,
tracking the trend in the Kramers escape time. The $\alpha=0$ to $\alpha=1/2$
ratio is $2.85$; this is close to the $2.77$ obtained from an isometric-coordinate
heuristic, but the nonlinear observation metric prevents treating that agreement
as a parameter-free prediction (Section~\ref{sec:res-spectra}).

\item \textbf{A metric that certifies its own failure.} At very small bandwidth
the graph becomes numerically nearly reducible and the leading eigenvector behaves
like a component indicator, yet $\lvert\rho_{\mathrm{Spearman}}\rvert$ reads $0.9997$ ---
locking onto residual variation eight orders of magnitude below the coordinate's
own range, because rank correlation discards magnitude by construction. We give a
MST-scale-constrained bandwidth selector and three diagnostics that detect
this (Sections~\ref{sec:bandwidth}, \ref{sec:degeneracy}).
\end{enumerate}

\subsection{Relation to a preliminary configuration}
\label{sec:pilot-intro}

An earlier configuration of this same project --- short horizon ($T=30$),
polynomial observations only, median-distance bandwidth, $\alpha = 1$ fixed ---
produced results that looked clean and were substantially misleading: it reported
$0.9846$ well-classification accuracy while measuring, as we now show, essentially
nothing about the wells. Four independent defects conspired, and each is
corrected here (Table~\ref{tab:corrections}). We document this explicitly in
Section~\ref{sec:disc-pilot}, because the failure mode is not specific to this
problem and each defect is individually plausible.

\begin{table}[t]
\centering
\caption{The preliminary configuration and the corrections made.}
\label{tab:corrections}
\small
\begin{tabular}{@{}p{0.29\textwidth}p{0.31\textwidth}p{0.32\textwidth}@{}}
\toprule
Defect & Consequence & Correction \\
\midrule
Horizon $T=30$, i.e.\ $2.4$ Kramers times; $2$--$3$ committed transitions;
$84.9\%$ occupancy of one well
& No dynamical statistics possible; the $2$-means threshold landed near $x=0$ by
luck
& $T = 600$ ($49$ escape times, $61$ transitions, $50.6\%$ occupancy) via time
subsampling at fixed $N$ \\[4pt]
Only $\alpha = 1$ run
& Density information removed; the decisive normalization comparison never made
& Sweep $\alpha \in \{0,\tfrac14,\tfrac12,\tfrac34,1\}$ \\[4pt]
Median-distance bandwidth
& Overshoots by $2\times$ (polynomial) to $29\times$ (RBF); drives the diffusion
map onto the PCA baseline
& MST-scale-constrained kernel-sum scan \\[4pt]
Polynomial observations only
& Manifold nearly straight, so PCA ties trivially and the comparison is
uninformative
& Add the RBF sensor-array model \\[4pt]
Pearson / Spearman / well accuracy only
& Cannot distinguish a monotone ramp from two compact levels, nor detect a
nearly reducible kernel
& Add well-alignment, two-cluster, and spectral degeneracy diagnostics \\
\bottomrule
\end{tabular}
\end{table}

\subsection{Outline}

Section~\ref{sec:background} reviews metastability and manifold learning.
Section~\ref{sec:model} specifies the generative model.
Section~\ref{sec:dm} derives the diffusion map, the $\alpha$-family, and
bandwidth selection. Section~\ref{sec:eval} defines the baseline and metrics,
Section~\ref{sec:implementation} the software, Section~\ref{sec:results} the
results, Section~\ref{sec:discussion} their interpretation,
Section~\ref{sec:limitations} limitations, and Section~\ref{sec:conclusion}
concludes.

% =============================================================================
\section{Background and Related Work}
\label{sec:background}
% =============================================================================

\subsection{Metastability and reaction coordinates}

For a reversible diffusion $\dd X_t = -\nabla V(X_t)\dd t + \sigma\dd W_t$ the
transition semigroup has generator
\begin{equation}
\mathcal{L} f = -\nabla V \cdot \nabla f + \tfrac{\sigma^2}{2}\Delta f ,
\label{eq:generator}
\end{equation}
self-adjoint in $L^2(\rho_s)$ with respect to the Boltzmann density
$\rho_s \propto e^{-2V/\sigma^2}$. When $V$ has deep minima separated by high
barriers, $\mathcal{L}$ has a cluster of eigenvalues near zero separated by a gap
from the rest; the associated eigenfunctions are nearly constant within each
basin and vary sharply across transition regions. These are the \emph{slow
modes}, and the leading nontrivial one is the \emph{reaction coordinate}
\citep{nadler2006diffusion}.

The point that organizes this paper is that this structure is a property of the
\emph{dynamics} --- specifically of the invariant density, whose bimodality is
what makes the transition rare. It is not a property of the geometry of the set
the trajectory lives on, which for a one-dimensional system is an interval.

\subsection{Diffusion maps}

Diffusion maps \citep{coifman2006diffusion}, building on Laplacian eigenmaps
\citep{belkin2003laplacian} and contemporaneous with Isomap
\citep{tenenbaum2000global} and LLE \citep{roweis2000nonlinear}, construct a
Markov chain on the data whose eigenvectors give a nonlinear embedding. Their
distinguishing feature is a one-parameter family of kernel normalizations,
indexed by $\alpha \in [0,1]$, controlling how much the sampling density
influences the resulting operator: $\alpha = 0$ gives the density-weighted graph
Laplacian, $\alpha = 1/2$ a backward Fokker--Planck-type operator in the observed
manifold metric, and $\alpha = 1$ the Laplace--Beltrami operator with density
removed. Convergence rates of the
graph operator to its continuum limit were established by
\citet{singer2006graph}.

Applications to stochastic dynamics were developed by \citet{nadler2006diffusion}
and \citet{coifman2008diffusion}, who showed that under the appropriate
normalization the eigenvectors approximate the eigenfunctions of the backward
Kolmogorov operator. Later refinements address the bandwidth: locally scaled maps
\citep{rohrdanz2011determination} and variable-bandwidth kernels
\citep{berry2016variable} adapt $\eps$ to local density, which matters greatly
when the density is as non-uniform as a metastable trajectory makes it.
\citet{banisch2017diffusion} construct kernels tailored to arbitrary
non-degenerate It\^o processes so that the operator approximates the true
generator rather than a density-biased proxy.

\subsection{Positioning}

This is a controlled, ground-truth-available study with a from-scratch
implementation. Its contribution is diagnostic: by comparing against an exactly
known latent state, against a linear baseline on identical data, and across two
observation manifolds of deliberately different curvature, it isolates the
conditions under which the standard construction does and does not deliver what
practitioners want from it.

% =============================================================================
\section{The Generative Model}
\label{sec:model}
% =============================================================================

\subsection{The double-well stochastic differential equation}

The hidden state evolves by overdamped Langevin dynamics in
\begin{equation}
V(x) = -\tfrac12 x^2 + \tfrac14 x^4 ,
\qquad -V'(x) = x - x^3 ,
\label{eq:potential}
\end{equation}
giving the scalar It\^o equation
\begin{equation}
\dd X_t = \bigl(X_t - X_t^3\bigr)\dd t + \sigma\dd W_t , \qquad X_0 = x_0 .
\label{eq:sde}
\end{equation}
The drift has stable fixed points at $x = \pm1$ (the wells) and an unstable one
at $x = 0$ (the barrier), with $b'(\pm1) = -2$ and $b'(0) = +1$. The barrier
height is $\Delta V = V(0) - V(\pm1) = 1/4$. Since the drift is a gradient the
process is reversible, with stationary density
\begin{equation}
\rho_s(x) = Z^{-1}\exp\!\left(-\frac{2V(x)}{\sigma^2}\right) ,
\label{eq:stationary}
\end{equation}
bimodal for the values of $\sigma$ used here.

\begin{remark}[The drift is not globally Lipschitz]
The cubic term makes $b$ only locally Lipschitz, satisfying a one-sided Lipschitz
condition. Standard strong-convergence theorems for explicit Euler--Maruyama
assume global Lipschitz continuity, and explicit Euler is known to diverge in the
strong sense for superlinearly growing drifts if the step is not controlled. In
practice the restoring cubic makes the scheme well behaved for small $\Delta t$;
the implementation checks finiteness of the whole trajectory and raises a
diagnostic error otherwise.
\end{remark}

\subsection{Kramers timescale and the design of the simulation}
\label{sec:kramers}

The escape rate from a well is given asymptotically by Kramers' formula
\citep{kramers1940brownian}; in the overdamped one-dimensional case with
$D = \sigma^2/2$,
\begin{equation}
k \;\approx\; \frac{\sqrt{V''(\pm 1)\,\lvert V''(0)\rvert}}{2\pi}
\exp\!\left(-\frac{\Delta V}{D}\right)
= \frac{\sqrt2}{2\pi}\exp\!\left(-\frac{1}{2\sigma^2}\right) ,
\label{eq:kramers}
\end{equation}
using $V''(\pm1) = 2$ and $\lvert V''(0)\rvert = 1$. At the default
$\sigma = 0.7$ this gives $\Delta V/D = 1.020$, $k \approx 0.0811$, and an escape
time $\tau_{\mathrm{esc}} = 1/k \approx 12.3$.

\begin{remark}[The barrier is shallow, deliberately]
$\Delta V/D \approx 1$ is not a deep barrier and Eq.~\eqref{eq:kramers}, an
asymptotic result for $\Delta V/D \gg 1$, is being used near the edge of its
validity. The consequence is visible in the data: the path recrosses the barrier
many times before committing, so raw sign changes badly overcount transitions
(Section~\ref{sec:res-trajectory}). We keep $\sigma = 0.7$ as the primary setting
anyway, because Section~\ref{sec:res-barrier} shows it is the \emph{most
discriminating} regime: when the barrier is deep, the point cloud becomes
geometrically clustered on its own and every choice of $\alpha$ succeeds, so the
distinction this paper is about would be invisible.
\end{remark}

\subsection{Euler--Maruyama discretization and time subsampling}
\label{sec:em}

Equation~\eqref{eq:sde} is integrated with the explicit Euler--Maruyama scheme
\citep{higham2001algorithmic,kloeden1992numerical}
\begin{equation}
X_{n+1} = X_n + \bigl(X_n - X_n^3\bigr)\Delta t + \sigma\sqrt{\Delta t}\,Z_n ,
\qquad Z_n \stackrel{\text{iid}}{\sim}\mathcal{N}(0,1) ,
\label{eq:em}
\end{equation}
obtained by replacing $\int_{t_n}^{t_{n+1}}\sigma\dd W_s$ with $\sigma\Delta W_n$
and using $\Delta W_n \sim \mathcal{N}(0,\Delta t)$ --- whence the
$\sqrt{\Delta t}$ scaling that distinguishes stochastic from deterministic
integrators. The test suite checks this variance scaling explicitly.

\begin{remark}[Order of convergence]
Euler--Maruyama has strong order $1/2$ and weak order $1$ in general. Here the
noise is \emph{additive}: $g(x)\equiv\sigma$ is constant, so $g'=0$, the Milstein
correction $\tfrac12 gg'(\Delta W^2 - \Delta t)$ vanishes identically,
Euler--Maruyama and Milstein coincide, and the scheme attains strong order $1$.
\end{remark}

Two requirements now conflict. Resolving the \emph{dynamics} demands a horizon of
many escape times, i.e.\ many integration steps; the diffusion map is $O(N^2)$ in
memory and time (Section~\ref{sec:complexity}) and so caps $N$ at a few thousand.
We resolve this by \textbf{subsampling in time}: the SDE is integrated at the fine
step $\Delta t = 0.01$, needed for accuracy and stability, but only every
$20$\textsuperscript{th} point is retained. The kept series spans
$T = 599.8 \approx 49\,\tau_{\mathrm{esc}}$ while holding $N = 3000$ points. This
is legitimate because the retained points still come from the same trajectory and
so are still distributed according to the invariant density --- which is the only
property the kernel construction uses. It costs temporal resolution of individual
crossings, which we do not need, and buys ergodic coverage, which we do. Where
transition statistics are reported they are computed on the \emph{unsubsampled}
path, since subsampling would miss short excursions.

\subsection{Two high-dimensional observation models}
\label{sec:obs}

The learning algorithms never see $X$. Each scalar state is mapped through a
nonlinear feature map $\phi:\R\to\R^m$, lifted by a fixed seeded Gaussian matrix
$A \in \R^{m\times d}$, corrupted, and standardized:
\begin{equation}
Y = F A + \eta E ,
\qquad F_{i\cdot} = \phi(X_i)^\top ,
\qquad E_{ik}\stackrel{\text{iid}}{\sim}\mathcal{N}(0,1) ,
\label{eq:observation}
\end{equation}
with $d = 15$ and $\eta = 0.1$ by default. Column standardization ensures
Euclidean distances --- on which the whole kernel construction rests --- are not
dominated by whichever coordinates carry the largest scale.

\paragraph{The polynomial map ($m=9$).}
\begin{equation}
\phi(x) = \bigl[\,x,\;x^2,\;x^3,\;\sin x,\;\cos x,\;\sin 2x,\;\cos 2x,\;
\tanh x,\;e^{-x^2}\,\bigr]^\top .
\label{eq:featpoly}
\end{equation}
This is the map used in the preliminary configuration. It looks aggressively
nonlinear, but four of its nine components are near-monotone in $x$, so its image
is a gently bent arc with a strong preferred direction.

\paragraph{The RBF sensor-array map ($m=12$).}
\begin{equation}
\phi_j(x) = \exp\!\left(-\frac{(x-c_j)^2}{2w^2}\right) ,
\qquad c_j \text{ uniform on } [-2.2,\,2.2],
\qquad w = 0.35 .
\label{eq:featrbf}
\end{equation}
Each channel is a localized detector responding only to a band of positions ---
the situation in spectrometers, place cells, and sensor arrays. States far apart
in $x$ excite disjoint channels and are nearly orthogonal, so the image curve is
strongly curved and spread across many dimensions. Its maximum-variance direction
need not order it monotonically. This is a regime in which a manifold method can
beat unsupervised PCA; it does not imply that every possible supervised linear
projection must fail. Including this map and the gently bent polynomial map makes
the PCA comparison informative.

\subsection{What is identifiable}
\label{sec:identifiability}

At $\eta = 0$ the observations $Y_{i\cdot} = A^\top\phi(X_i)$ are an injective
deterministic function of $X_i$ alone,\footnote{Injective on the trajectory's
range: for the polynomial map the component $x\mapsto x$ is injective; for the
RBF map the vector of responses determines $x$ within the covered range. $A$ is
almost surely of full row rank $m \le d$.} so the noiseless cloud lies exactly on
a one-dimensional curve $\mathcal{M}\subset\R^{15}$ parameterized by $X$. The
intrinsic dimension is $1$; the ambient dimension is $15$.

No unsupervised method can recover $X$ itself, because nothing distinguishes $X$
from any strictly monotone reparameterization $h(X)$ --- the point cloud is
identical. \emph{The strongest achievable result is recovery of $X$ up to a
monotone reparameterization and a global sign.} This is why Spearman rank
correlation, not Pearson, is the fair headline metric for Q1: it is invariant
under $h$, whereas Pearson penalizes a warping that is fundamentally
unidentifiable. We report both, reading Spearman as recovery and the
Pearson--Spearman gap as a measure of warping.

% =============================================================================
\section{Diffusion Maps}
\label{sec:dm}
% =============================================================================

Let $y_1,\dots,y_N\in\R^d$ be the rows of the standardized observation matrix.

\subsection{Kernel and affinity graph}

\begin{equation}
D^2_{ij} = \norm{y_i-y_j}^2 ,
\qquad
K_{ij} = \exp\!\left(-\frac{D^2_{ij}}{\eps}\right) ,
\label{eq:kernel}
\end{equation}
with $\eps>0$ a \emph{squared} length scale. $K$ is symmetric, strictly positive,
and positive semidefinite, defining a weighted complete graph in which weight
decays on the scale $\sqrt\eps$.

\subsection{Anisotropic ($\alpha$) normalization and its generator limits}
\label{sec:alpha}

The data are sampled not uniformly from $\mathcal{M}$ but from the trajectory,
hence from the bimodal density~\eqref{eq:stationary}. Left uncorrected the graph
operator confounds geometry with density. \citet{coifman2006diffusion} remove the
confound with a preliminary normalization: writing $q_i = \sum_j K_{ij}$,
\begin{equation}
\Kt_{ij} = \frac{K_{ij}}{q_i^\alpha q_j^\alpha} ,
\qquad \alpha\in[0,1] .
\label{eq:alpha}
\end{equation}
As $N\to\infty$, $\eps\to0$ the associated operator converges, up to a positive
constant, to
\begin{equation}
\mathcal{L}_\alpha = \Delta
\;+\; 2\,\frac{\nabla\rho^{\,1-\alpha}}{\rho^{\,1-\alpha}}\cdot\nabla ,
\label{eq:alphafamily}
\end{equation}
with $\rho$ the sampling density \emph{with respect to volume on the observed
manifold} and $\Delta$ the Laplace--Beltrami operator for its induced metric.
This qualification matters: after a nonlinear observation map, $\rho$ is not
simply the state-space density $\rho_s(x)$, and $\Delta$ is not the ordinary
second derivative in $x$. The three canonical choices are therefore interpreted
in observation coordinates:

\begin{itemize}
\item $\alpha=0$: $\mathcal{L}_0 = \Delta + 2(\nabla\rho/\rho)\cdot\nabla$, the
classical normalized graph Laplacian. It emphasizes the observed sampling
density most strongly.
\item $\alpha=1/2$: $\mathcal{L}_{1/2} = \Delta + (\nabla\rho/\rho)\cdot\nabla$,
a backward Fokker--Planck-type operator that retains density information. If the
observations were an isometric coordinate for $x$, it would be a constant
multiple of the true generator~\eqref{eq:generator}. For the nonlinear maps used
here it is instead a dynamics-informed proxy in the induced observation metric.
\item $\alpha=1$: the drift term vanishes and $\mathcal{L}_1 = \Delta$, the pure
Laplace--Beltrami operator. The density is divided out completely, so the
eigenfunctions describe the \emph{shape} of the data set and nothing about the
dynamics that generated it.
\end{itemize}

In an isometric state coordinate, the $\alpha=0$ drift would correspond to a
doubled effective potential. Kramers' formula would then suggest that
$\mu_2/\mu_1$ is larger at $\alpha=0$ than at $\alpha=1/2$ by
$e^{\Delta V/D}=2.77$ at $\sigma=0.7$. We retain this as a heuristic benchmark in
Section~\ref{sec:res-spectra}, not as an exact prediction for the non-isometric
observation maps.

\subsection{Markov matrix and symmetric conjugation}

With degrees $d_i = \sum_j\Kt_{ij}$ and $D = \diag(d_i)$,
\begin{equation}
P = D^{-1}\Kt , \qquad \sum_j P_{ij} = 1 ,
\label{eq:markov}
\end{equation}
the transition matrix of a random walk on the data. In exact arithmetic all
entries are positive, so the chain is irreducible and aperiodic. At very small
bandwidth, however, floating-point weights can be negligible and the matrix can
be numerically nearly reducible, producing additional eigenvalues extremely close
to $1$. Away from that regime, $\lambda_0 = 1$ is simple with constant right
eigenvector and $\lvert\lambda_k\rvert<1$ otherwise. $P$ is
not symmetric, and a direct nonsymmetric eigensolve is slower and vulnerable to
spurious complex eigenvalues from floating-point asymmetry. The following removes
the difficulty.

\begin{proposition}[Symmetric conjugacy]
\label{prop:conjugacy}
Let $\Kt$ be symmetric with strictly positive degrees, $P = D^{-1}\Kt$, and
$P_{\mathrm{sym}} = D^{-1/2}\Kt D^{-1/2}$. Then $P_{\mathrm{sym}}$ is symmetric,
$P = D^{-1/2}P_{\mathrm{sym}}D^{1/2}$, so $P$ and $P_{\mathrm{sym}}$ share the
same real spectrum; and if $P_{\mathrm{sym}}v_k = \lambda_k v_k$ then
$\psi_k = D^{-1/2}v_k$ is a right eigenvector of $P$ with eigenvalue $\lambda_k$.
\end{proposition}

\begin{proof}
Symmetry is immediate from that of $\Kt$ and the diagonality of $D$. For the
conjugacy,
\[
D^{-1/2}P_{\mathrm{sym}}D^{1/2}
= D^{-1/2}\bigl(D^{-1/2}\Kt D^{-1/2}\bigr)D^{1/2}
= D^{-1}\Kt = P .
\]
Similar matrices share eigenvalues, and a real symmetric matrix has real ones.
Finally
\[
P\bigl(D^{-1/2}v_k\bigr)
= D^{-1/2}P_{\mathrm{sym}}D^{1/2}D^{-1/2}v_k
= D^{-1/2}P_{\mathrm{sym}}v_k
= \lambda_k D^{-1/2}v_k . \qedhere
\]
\end{proof}

The whole spectral computation is therefore performed on the real symmetric
positive semidefinite $P_{\mathrm{sym}}$ with \texttt{scipy.linalg.eigh}, which
guarantees real eigenvalues and an orthonormal basis and allows requesting only
the top eigenpairs. The implementation additionally symmetrizes
$P_{\mathrm{sym}}\leftarrow\tfrac12(P_{\mathrm{sym}}+P_{\mathrm{sym}}^\top)$
against round-off before calling \texttt{eigh}.

\subsection{Diffusion coordinates and diffusion distance}

Ordering $1 = \lambda_0 > \lambda_1 \ge \lambda_2 \ge \cdots$ and discarding the
trivial constant $\psi_0$, the diffusion map at time $t$ is
\begin{equation}
\Psi_t(y_i) = \bigl(\lambda_1^t\psi_1(i),\dots,\lambda_s^t\psi_s(i)\bigr) ,
\label{eq:embedding}
\end{equation}
whose Euclidean distance equals the \emph{diffusion distance}
\begin{equation}
D_t^2(y_i,y_j) = \sum_{k\ge1}\lambda_k^{2t}\bigl(\psi_k(i)-\psi_k(j)\bigr)^2
= \sum_\ell \frac{\bigl(P^t_{i\ell}-P^t_{j\ell}\bigr)^2}{\pi_\ell} ,
\label{eq:diffdist}
\end{equation}
with $\pi$ the stationary distribution. Two points are close when random walks
from each reach the rest of the data similarly after $t$ steps --- a notion
robust to noise because it aggregates over all paths. We use $t=1$, $s=5$.

\subsection{Bandwidth selection}
\label{sec:bandwidth}

The bandwidth is the most consequential hyperparameter, and it is bounded from
both sides. If $\eps$ is too large, $K \approx \mathbf{1}\mathbf{1}^\top$: the
walk mixes in one step and geometry is lost. If $\eps$ is too small,
$K\approx I$: the kernel becomes numerically nearly reducible and the leading
eigenvectors behave like component indicators.

\paragraph{The median heuristic and why it fails.}
The preliminary configuration set $\eps$ to the median of the nonzero pairwise
squared distances. This is scale-free and robust, but it has no connection to the
geometry of the data, and it overshoots: measured against the selector below it
is $2.0\times$ too large on the polynomial model and $29\times$ too large on the
RBF model. Section~\ref{sec:res-bandwidth} shows what that costs.

\paragraph{MST-scale-constrained kernel-sum scan.}
We instead use the criterion of \citet{coifman2006diffusion}. With
\begin{equation}
S(\eps) = \sum_{i,j}\exp\!\left(-\frac{D^2_{ij}}{\eps}\right) ,
\end{equation}
$S\to N$ as $\eps\to0$ and $S\to N^2$ as $\eps\to\infty$; between the plateaus
$\log S$ is approximately linear in $\log\eps$ with slope $m\approx
d_{\mathrm{intrinsic}}/2$, and that linear region is exactly the range at which
the kernel resolves the manifold. We take the bandwidth of maximum slope,
restricted from below by the conservative scale
\begin{equation}
\eps_{\mathrm{MST}} = \max\{\text{edge weights of the MST of } D^2\} ,
\label{eq:floor}
\end{equation}
which makes every spanning-tree edge have affinity at least $e^{-1}$. This is not
a literal connectivity threshold: a Gaussian graph is complete in exact
arithmetic, and any operational threshold depends on an explicit affinity
cutoff. It is a deliberately conservative lower scale that keeps the scan away
from the nearly reducible regime where metrics fail silently
(Section~\ref{sec:degeneracy}). The recovered slopes give
intrinsic-dimension estimates $2m = 1.13$ (polynomial) and $2.03$ (RBF), against
a true value of $1$ --- a useful sanity check, and a first hint that the RBF
manifold is the harder object.

\paragraph{An honest caveat.}
This selector is a large improvement but is not reliable. Its criterion assumes a
single manifold scale, whereas metastable data have two --- the within-well
spread and the between-well separation --- and the kernel-sum curve accordingly
has two knees. When the second dominates, the rule selects a bandwidth an order
of magnitude too large; re-selecting at every measurement-noise level made the
chosen $\eps$ jump between $0.63$ and $16.7$ on the RBF model. We therefore hold
$\eps$ fixed within the noise study (Section~\ref{sec:res-noise}) and treat the
bandwidth sweep, not the selector, as the primary evidence. A variable-bandwidth
kernel \citep{berry2016variable,rohrdanz2011determination} is the principled fix
and is left to future work.

\subsection{Computational complexity}
\label{sec:complexity}

The dense distance and kernel matrices are $N\times N$, so memory and time scale
as $O(N^2)$ and a full eigendecomposition as $O(N^3)$ (reduced here by requesting
only the top $s+1$ pairs). A \texttt{float64} $N\times N$ matrix occupies $8N^2$
bytes --- about $72$\,MB at $N=3000$, $800$\,MB at $N=10^4$ --- and several are
live at once; the implementation warns above $N=4000$. This wall is exactly what
the time subsampling of Section~\ref{sec:em} is designed to work around. Sparse
$k$-nearest-neighbour kernels and the Nystr\"om extension
\citep{fowlkes2004spectral} are the standard remedies and are out of scope here.
The full experiment --- four studies, $52$ embeddings, $17$ figures --- runs in a
few minutes.

\begin{algorithm}[t]
\caption{Diffusion map with anisotropic normalization}
\label{alg:dm}
\begin{algorithmic}[1]
\Require $Y\in\R^{N\times d}$ (standardized); $\alpha$; $s$; diffusion time $t$;
         bandwidth $\eps$ or the rule \texttt{scan}
\Ensure  coordinates $\Psi\in\R^{N\times s}$, eigenvalues, diagnostics
\State $D^2_{ij}\gets\norm{y_i-y_j}^2$
       \Comment{via \texttt{pdist}/\texttt{squareform}}
\If{$\eps = \texttt{scan}$}
  \State $\eps_{\mathrm{MST}}\gets$ largest MST edge of $D^2$
         \Comment{conservative scale, Eq.~\eqref{eq:floor}}
  \State $\eps\gets\arg\max_{\eps\ge\eps_{\mathrm{MST}}}
         \dd\log S(\eps)/\dd\log\eps$
\EndIf
\State $K_{ij}\gets\exp(-D^2_{ij}/\eps)$
\State $q_i\gets\sum_j K_{ij}$;\quad
       $\Kt_{ij}\gets K_{ij}/(q_i^\alpha q_j^\alpha)$
\State $d_i\gets\sum_j\Kt_{ij}$;\quad \textbf{assert} $d_i>0$
\State $P_{\mathrm{sym}}\gets D^{-1/2}\Kt D^{-1/2}$;\quad
       $P_{\mathrm{sym}}\gets\tfrac12(P_{\mathrm{sym}}+P_{\mathrm{sym}}^\top)$
\State $(\lambda_k,v_k)_{k=0}^s\gets$ top $s{+}1$ eigenpairs of
       $P_{\mathrm{sym}}$ \Comment{\texttt{eigh}}
\State $\psi_k\gets D^{-1/2}v_k$ \Comment{Proposition~\ref{prop:conjugacy}}
\State $n_{\mathrm{triv}}\gets\#\{k:\lambda_k>1-10^{-8}\}$
       \Comment{$>1 \Rightarrow$ disconnected or nearly reducible; warn}
\State discard $(\lambda_0,\psi_0)$; \quad
       $\Psi_{\cdot k}\gets\max(\lambda_k,0)^t\,\psi_k$
\State \Return $\Psi$, $(\lambda_k)$, $n_{\mathrm{triv}}$,
       $\#\{\text{distinct values of }\Psi_{\cdot1}\}$
\end{algorithmic}
\end{algorithm}

% =============================================================================
\section{Baseline and Evaluation Protocol}
\label{sec:eval}
% =============================================================================

\subsection{PCA baseline}

PCA is applied to the \emph{identical} standardized $Y$, so the methods differ
only in the embedding. PC$_1$ is the best linear one-dimensional summary;
$\psi_1$ a nonlinear one. If $\mathcal{M}$ is strongly curved in $\R^{15}$ a
single linear direction cannot follow it and PCA should lose ground; if
$\mathcal{M}$ is nearly straight the two should agree. With two observation
models of deliberately different curvature, the comparison becomes a measurement
of that prediction rather than a single data point.

\subsection{Metrics}

Every coordinate has an arbitrary global sign and an unidentifiable monotone
warping (Section~\ref{sec:identifiability}). All reported metrics are made
sign-invariant, but only Spearman correlation is invariant to arbitrary monotone
warping. The first three were used in the preliminary configuration; the last three are
added here, and Section~\ref{sec:disc-pilot} shows why they were necessary.

\begin{description}
\item[$\lvert r\rvert$, Pearson with $X$.] Linear association. Sensitive to
monotone warping, so informative about \emph{how much} a method warps rather than
as a pass/fail test.

\item[$\lvert\rho\rvert$, Spearman with $X$.] Monotone association: the headline
metric for Q1, exactly invariant to the unidentifiable reparameterization.

\item[Well accuracy.] Balanced accuracy of a $2$-means split of the coordinate
against true membership $\mathbf{1}\{X\ge0\}$, with cluster labels matched in
whichever orientation scores higher. Balanced accuracy is used because the wells
need not be equally occupied; chance is $0.5$.

\item[Well-indicator alignment, $\lvert r(\psi_1,\sgn X)\rvert$.] Association
with well membership rather than position. A coordinate behaving like a smoothed
well indicator scores highly. This is a supervised shape diagnostic, not by
itself a measurement of a dynamical timescale.

\item[Two-cluster score.] Variance fraction explained by a sample-weighted
$2$-means partition. Values near $1$ indicate two compact
levels, but the score is not a formal test of bimodality: even a uniform interval
scores $0.75$. We therefore use it only descriptively, relative to the same score
on the true state ($0.829$).

\item[Degeneracy diagnostics.] The number $n_{\mathrm{triv}}$ of computed
eigenvalues within $10^{-8}$ of $1$ and the fraction of distinct values taken by
$\psi_1$. Exact components create repeated unit eigenvalues; weak links create
near-unit modes, so $n_{\mathrm{triv}}$ is not interpreted as a literal component
count.
\end{description}

\begin{remark}[Why well accuracy alone is not enough]
\label{rem:wellscore}
A $2$-means split of \emph{any} strictly monotone function of $X$ places a
threshold somewhere in the range of $X$; if it happens to fall near $x=0$ the
balanced accuracy is high. A high well score is therefore necessary but not
sufficient evidence of recovered two-state structure --- genuine separation
should additionally be visible in the coordinate distribution and its
two-cluster variance score. In the preliminary
configuration the wells were $85/15$ occupied and the threshold landed near zero
by luck, producing $0.9846$ from a coordinate with no two-state content.
\end{remark}

\subsection{Studies}

Four sweeps, each varying one factor: \textbf{$\alpha$} over
$\{0,\tfrac14,\tfrac12,\tfrac34,1\}$ at each model's selected bandwidth;
\textbf{bandwidth} over $\{2^{-2},\dots,2^5\}\times\eps_{\mathrm{scan}}$;
\textbf{measurement noise} over $\{0,0.05,0.1,0.2,0.4\}$ at fixed $\eps$; and
\textbf{barrier depth} over $\sigma\in\{0.45,0.5,0.55,0.7\}$ with the horizon
held fixed in units of the Kramers time, so that every setting sees a comparable
number of transitions and depth is not confounded with sample size.

% =============================================================================
\section{Implementation and Reproducibility}
\label{sec:implementation}
% =============================================================================

Each stage occupies a focused module. \texttt{src/simulation.py} handles the SDE
and transition statistics; \texttt{src/observations.py} handles both observation
maps; and \texttt{src/diffusion\_maps.py} implements Algorithm~\ref{alg:dm}.
Metrics and PCA are in \texttt{src/evaluation.py}. The studies, plotting, and CLI
are in \texttt{src/experiments.py}, \texttt{src/visualization.py}, and
\texttt{run\_experiment.py}, respectively. Diffusion maps are implemented directly from the
mathematics of Section~\ref{sec:dm}; no manifold-learning library is used. NumPy
\citep{harris2020array}, SciPy \citep{virtanen2020scipy}, scikit-learn
\citep{pedregosa2011scikit} (PCA, $2$-means, balanced accuracy only), pandas, and
Matplotlib \citep{hunter2007matplotlib} supply primitives.

Configuration is layered: dataclass defaults, then \texttt{config.yaml}, then
command-line flags. Each stage is configured by a frozen dataclass validating its
own parameters on construction, and numerical failure modes are trapped where
they occur --- non-finite trajectories, non-positive graph degrees, and
disconnected or nearly reducible kernels. All randomness flows through explicitly seeded
\texttt{numpy.random.Generator} objects, and the resolved parameter set is written
to \texttt{outputs/metrics/run\_parameters.json} on every run. The complete
experiment is reproduced by \texttt{python run\_experiment.py} with no network
access and no external data.

A \texttt{pytest} suite of $55$ tests passes in full. The
tests are properties, not golden values: bit-identical reproduction under a fixed
seed, the $\sqrt{\Delta t}$ noise scaling, row-stochasticity of $P$,
$\lambda_0 = 1$ with constant eigenvector, agreement of the symmetric-conjugation
eigenvectors with a direct nonsymmetric solve, and --- added for this revision ---
that subsampling returns a strict subset of the fine path, that committed
transitions ignore barrier chatter, that the RBF image is less linearly
explainable than the polynomial one, that the scan respects its conservative MST
scale, that a small-bandwidth near-unit spectrum is flagged degenerate, that the
two-cluster score is not mistaken for a modality test, and that the gap-ratio
helper reproduces the $k^2$ law exactly on synthetic input.

% =============================================================================
\section{Results}
\label{sec:results}
% =============================================================================

\begin{table}[t]
\centering
\caption{Resolved run parameters
(\texttt{outputs/metrics/run\_parameters.json}).}
\label{tab:params}
\small
\begin{tabular}{@{}llr@{}}
\toprule
Stage & Parameter & Value \\
\midrule
Simulation   & retained samples $N$          & $3000$ \\
             & integration step $\Delta t$   & $0.01$ \\
             & subsample stride              & $20$ \\
             & horizon $T$                   & $599.8$ \\
             & $T/\tau_{\mathrm{esc}}$       & $48.7$ \\
             & process noise $\sigma$        & $0.7$ \\
             & seed                          & $42$ \\
\midrule
Observation  & ambient dimension $d$         & $15$ \\
             & polynomial features $m$       & $9$ \\
             & RBF channels / width          & $12$ / $0.35$ \\
             & measurement noise $\eta$      & $0.1$ \\
             & observation seed              & $123$ \\
\midrule
Diffusion map& bandwidth rule                & scan \\
             & normalization $\alpha$        & $0.5$ \\
             & components $s$                & $5$ \\
             & diffusion time $t$            & $1.0$ \\
\bottomrule
\end{tabular}
\end{table}

\subsection{The latent trajectory and sampling adequacy}
\label{sec:res-trajectory}

Figure~\ref{fig:traj} shows the simulated state and its density: long dwells near
$x=\pm1$ punctuated by abrupt crossings, and a clearly bimodal density matching
Eq.~\eqref{eq:stationary}.

\begin{figure}[t]
\centering
\begin{subfigure}[b]{0.62\textwidth}
  \includegraphics[width=\linewidth]{01_state_vs_time.png}
  \caption{Trajectory $X_t$ over $T = 600$.}
\end{subfigure}
\hfill
\begin{subfigure}[b]{0.36\textwidth}
  \includegraphics[width=\linewidth]{02_state_histogram.png}
  \caption{Empirical density.}
\end{subfigure}
\caption{The hidden state. Occupancy is now nearly balanced ($50.6\%$ right
well) across $61$ committed transitions, against $15.1\%$ and $2$--$3$
transitions in the preliminary configuration.}
\label{fig:traj}
\end{figure}

Quantitatively the sampling is now adequate for dynamical statistics. Over
$T = 599.8 = 48.7\,\tau_{\mathrm{esc}}$ the fine path records $587$ sign changes
but only $61$ \emph{committed} transitions --- those entering the core of one
well ($\lvert x\rvert>0.5$) having last been in the core of the other. The gap
between $587$ and $61$ is barrier chatter, the expected consequence of
$\Delta V/D \approx 1.02$, and is exactly why a raw sign-change count is the
wrong statistic. Against the Kramers prediction of $48.7$ escapes the measured
$61$ is the right order, with the excess consistent with
Eq.~\eqref{eq:kramers} underestimating the rate for a shallow barrier. Right-well
occupancy is $0.506$ and the two-cluster variance score of $X$ is $0.829$, the
descriptive reference used below.

\subsection{The two observation manifolds}
\label{sec:res-manifolds}

Figure~\ref{fig:features} shows feature traces for both models. PCA on the
standardized observations gives explained-variance ratios
\begin{align}
\text{polynomial:}\quad &(0.641,\ 0.192,\ 0.147,\ 0.015,\ 0.003),
\quad \textstyle\sum_{1}^{3} = 0.980 , \\
\text{RBF:}\quad &(0.519,\ 0.241,\ 0.128,\ 0.068,\ 0.016),
\quad \textstyle\sum_{1}^{3} = 0.888 .
\end{align}
Both are effectively low-dimensional in $\R^{15}$, confirming the random lift
added ambient dimension without information. But the leading component tells the
story: $64\%$ on one linear direction for the polynomial map means a gently bent
arc; $52\%$, with a much slower decay, means the RBF image genuinely resists
linear summary. The intrinsic-dimension estimates from the bandwidth scan agree
($1.13$ versus $2.03$ against a true value of $1$; the RBF overestimate reflects
that at the selected scale the kernel still sees the noise directions).

\begin{figure}[t]
\centering
\begin{subfigure}[b]{0.49\textwidth}
  \includegraphics[width=\linewidth]{03_poly_feature_traces.png}
  \caption{Polynomial features.}
\end{subfigure}
\hfill
\begin{subfigure}[b]{0.49\textwidth}
  \includegraphics[width=\linewidth]{03_rbf_feature_traces.png}
  \caption{RBF sensor channels.}
\end{subfigure}
\caption{Observation features. Individually uninformative; collectively an
injective encoding of $X$. The RBF channels fire selectively as the state passes
through each detector's band.}
\label{fig:features}
\end{figure}

\subsection{Q1: state recovery, and where PCA fails}
\label{sec:res-q1}

Table~\ref{tab:headline} and Figure~\ref{fig:recovery} give the central result
for Q1.

\begin{table}[t]
\centering
\caption{Headline comparison at the default settings. The reference two-cluster
variance score of the true state is $0.829$.}
\label{tab:headline}
\begin{tabular}{@{}llrrrr@{}}
\toprule
Observations & Method & $\lvert\rho\rvert$ & well alignment & 2-cluster score
& well acc. \\
\midrule
polynomial & PCA                 & $0.9997$ & $0.8760$ & $0.768$ & $0.9852$ \\
           & DM $\alpha=1/2$     & $0.9998$ & $0.8407$ & $0.707$ & $0.9937$ \\
           & DM $\alpha=1$       & $0.9998$ & $0.6670$ & $0.511$ & $0.5804$ \\
\midrule
RBF        & PCA                 & $0.6624$ & $0.9309$ & $0.878$ & $0.9795$ \\
           & DM $\alpha=1/2$     & $0.9998$ & $0.9644$ & $0.933$ & $0.9885$ \\
           & DM $\alpha=1$       & $0.9998$ & $0.9407$ & $0.885$ & $0.9957$ \\
\bottomrule
\end{tabular}
\end{table}

\begin{figure}[t]
\centering
\begin{subfigure}[b]{0.245\textwidth}
  \includegraphics[width=\linewidth]{04_poly_dm_coordinate_vs_state.png}
  \caption{poly, DM}
\end{subfigure}
\begin{subfigure}[b]{0.245\textwidth}
  \includegraphics[width=\linewidth]{05_poly_pca_coordinate_vs_state.png}
  \caption{poly, PCA}
\end{subfigure}
\begin{subfigure}[b]{0.245\textwidth}
  \includegraphics[width=\linewidth]{04_rbf_dm_coordinate_vs_state.png}
  \caption{RBF, DM}
\end{subfigure}
\begin{subfigure}[b]{0.245\textwidth}
  \includegraphics[width=\linewidth]{05_rbf_pca_coordinate_vs_state.png}
  \caption{RBF, PCA}
\end{subfigure}
\caption{The leading coordinate against the true state. Three panels show a tight
monotone curve; panel (d) does not. On the RBF observations PC$_1$ rises to a
peak near $X\approx0.5$ and falls thereafter, so it assigns the same value to
$X\approx-0.3$ and $X\approx1.5$. Thus the maximum-variance linear projection
chosen by PCA does not order this manifold, and its rank correlation drops to
$0.6624$.}
\label{fig:recovery}
\end{figure}

On the polynomial observations the two methods tie: $0.9997$ against $0.9998$,
with PCA's Pearson correlation ($0.9932$) even slightly \emph{ahead} of the
diffusion map's at $\alpha=1/2$ ($0.9554$), since the diffusion coordinate applies
a mild monotone warp that Pearson penalizes and Spearman correctly ignores. A
study run only on this model --- the preliminary configuration --- would conclude
that manifold learning buys nothing.

On the RBF observations the comparison separates decisively:
$\lvert\rho\rvert = 0.9998$ for the diffusion map against $0.6624$ for PCA.
Figure~\ref{fig:recovery}(d) shows the mechanism directly: PC$_1$ is a
non-monotone hump. The maximum-variance linear summary folds this curve, and
folding destroys the ordering. Figure~\ref{fig:embeddings} makes the same point
in two dimensions.

\begin{figure}[t]
\centering
\begin{subfigure}[b]{0.48\textwidth}
  \includegraphics[width=\linewidth]{06_rbf_dm_embedding_scatter.png}
  \caption{Diffusion map.}
\end{subfigure}
\hfill
\begin{subfigure}[b]{0.48\textwidth}
  \includegraphics[width=\linewidth]{07_rbf_pca_embedding_scatter.png}
  \caption{PCA.}
\end{subfigure}
\caption{Two-dimensional embeddings of the RBF observations, coloured by the true
state.}
\label{fig:embeddings}
\end{figure}

The answer to Q1 is therefore: \emph{yes, and curvature strongly affects the
margin}. On these two encodings the diffusion map ties PCA on the gently bent
curve and does dramatically better on the strongly curved one. This comparison
concerns unsupervised PCA, not every possible supervised linear projection.

\subsection{Q2: normalization and observation geometry}
\label{sec:res-alpha}

Table~\ref{tab:alpha} and Figure~\ref{fig:alpha} sweep $\alpha$ with everything
else held fixed.

\begin{table}[t]
\centering
\caption{The $\alpha$ study. State recovery is flat. Density-sensitive shape
measures fall strongly on the polynomial model and only weakly on the
geometrically separated RBF model. Reference two-cluster score of $X$: $0.829$.}
\label{tab:alpha}
\begin{tabular}{@{}llrrrrr@{}}
\toprule
Observations & $\alpha$ & $\lvert\rho\rvert$ & well alignment & 2-cluster score
& well acc. & $\mu_2/\mu_1$ \\
\midrule
polynomial & $0.00$ & $0.9998$ & $0.9071$ & $0.823$ & $0.9960$ & $2.33$ \\
           & $0.25$ & $0.9998$ & $0.8924$ & $0.797$ & $0.9960$ & $1.90$ \\
           & $0.50$ & $0.9998$ & $0.8407$ & $0.707$ & $0.9937$ & $1.41$ \\
           & $0.75$ & $0.9998$ & $0.6566$ & $0.446$ & $0.9249$ & $1.58$ \\
           & $1.00$ & $0.9998$ & $0.6670$ & $0.511$ & $\mathbf{0.5804}$ & $2.99$ \\
\midrule
RBF        & $0.00$ & $0.9997$ & $0.9686$ & $0.949$ & $0.9754$ & $28.99$ \\
           & $0.25$ & $0.9997$ & $0.9675$ & $0.942$ & $0.9825$ & $16.92$ \\
           & $0.50$ & $0.9998$ & $0.9644$ & $0.933$ & $0.9885$ & $10.17$ \\
           & $0.75$ & $0.9998$ & $0.9571$ & $0.916$ & $0.9939$ & $6.27$ \\
           & $1.00$ & $0.9998$ & $0.9407$ & $0.885$ & $0.9957$ & $4.19$ \\
\bottomrule
\end{tabular}
\end{table}

\begin{figure}[t]
\centering
\includegraphics[width=\textwidth]{08_alpha_study.png}
\caption{The normalization study. Left: state recovery is untouched by $\alpha$
(and the RBF PCA baseline sits far below). Centre and right: density-sensitive
shape changes strongly on the polynomial model and weakly on the RBF model.}
\label{fig:alpha}
\end{figure}

Three things are visible at once. First, \textbf{$\alpha$ does not affect Q1}:
$\lvert\rho\rvert$ is $0.9997$--$0.9998$ everywhere. Geometric recovery is
robust to the normalization.

Second, \textbf{$\alpha$ strongly controls the polynomial-model answer}. Well
alignment falls from $0.907$ to $0.667$, the two-cluster score from $0.823$ to
$0.511$ against a true-state reference of $0.829$, and well accuracy from
$0.9960$ to $\mathbf{0.5804}$ --- poor, though above the $0.5$ chance baseline.
The data, bandwidth, and eigensolver are fixed; only the density normalization
changes.

Third, \textbf{the two-cluster score gives a descriptive view of the same trend}.
At $\alpha=0$ the polynomial model gives $0.823$ against the true-state reference
$0.829$; at $\alpha=1/2$ it is $0.707$ and at $\alpha=1$, $0.511$. These values
show loss of compact two-level structure as density is removed, but they are not
a formal test of bimodality and should not be read as exact operator theory.

The RBF model shows the same well-alignment trend but far more weakly
$0.969\to0.941$; well accuracy actually \emph{rises}). The reason is instructive
and is confirmed independently in Section~\ref{sec:res-barrier}: in the RBF
encoding, states in different wells excite disjoint channels and are therefore
far apart \emph{geometrically}, so the point cloud is already two-clustered as a
set. Geometry and dynamics happen to agree, and dividing out the density costs
little. On the polynomial model, whose image is a compact arc, density supplies
most of the visible well separation, and $\alpha=1$ suppresses it.

\subsection{Spectral signatures: which operator was computed}
\label{sec:res-spectra}

The eigenvalues identify the limiting operator, and do so more sharply than any
of the correlation metrics. We read them through the gap ratios
$(1-\lambda_k)/(1-\lambda_1)$, which estimate $\mu_k/\mu_1$ since
$\lambda_k \approx 1 - c\,\eps\,\mu_k$ and the unknown constant cancels.

\paragraph{At $\alpha=1$: the $k^2$ law of an interval.}
For the Laplace--Beltrami operator on a one-dimensional interval with free ends
the eigenfunctions are the Neumann cosine modes $\cos(k\pi s/L)$ with
$\mu_k = (k\pi/L)^2$, so $\mu_k/\mu_1 = k^2$, i.e.\ $(4,\,9,\,16)$. The RBF model
at $\alpha=1$ measures $(4.19,\ 8.38)$ for $k=2,3$. The diffusion map has computed
the Laplacian of a curve.

\paragraph{At $\alpha=1/2$: a density-sensitive timescale gap.}
The same model at $\alpha=1/2$ gives $\mu_2/\mu_1 = 10.17$ --- a genuine
separation between one slow mode and the rest, which is the spectral signature of
metastable structure in the sampled cloud and is nothing like $4$. Because the
RBF observation map is not isometric, this is not an exact physical timescale of
Eq.~\eqref{eq:sde}.

\paragraph{A heuristic check on $\alpha=0$.}
In an isometric state coordinate, Section~\ref{sec:alpha} gives the heuristic
factor $e^{\Delta V/D}=e^{1.020}=2.77$ between the $\alpha=0$ and $\alpha=1/2$
gap ratios. The measured ratio is $28.99/10.17=\mathbf{2.85}$. The proximity is
consistent with stronger density weighting at $\alpha=0$, but it is not a
parameter-free validation: Kramers' formula is being used outside its deep-barrier
regime, and the induced RBF metric modifies both operators.

\begin{figure}[t]
\centering
\includegraphics[width=0.62\textwidth]{12_alpha_spectra.png}
\caption{Timescale separation versus normalization. The RBF curve descends from
$29.0$ at $\alpha=0$ to $4.19$ at $\alpha=1$ --- landing on the dashed $k^2$ line
that marks the one-dimensional interval Laplacian.}
\label{fig:spectra}
\end{figure}

The polynomial model at $\alpha=1$ gives $(2.99,\ 5.75)$, heading toward
$(4,\,9)$ but not arriving, because its conservative MST scale leads to a bandwidth
($\eps = 12.8$) too large for the continuum limit to be well approximated. This
is itself a useful observation: the $k^2$ signature is clearest exactly when the
bandwidth can be taken small, which is when the manifold is well spread out.

\subsection{Sensitivity to kernel bandwidth}
\label{sec:res-bandwidth}

\begin{table}[t]
\centering
\caption{Bandwidth study at $\alpha=1/2$. PCA is bandwidth-independent
($0.9997$ polynomial, $0.6624$ RBF).}
\label{tab:bandwidth}
\begin{tabular}{@{}lrrrr@{}}
\toprule
$\eps/\eps_{\mathrm{scan}}$ & \multicolumn{2}{c}{polynomial}
& \multicolumn{2}{c}{RBF} \\
\cmidrule(lr){2-3}\cmidrule(lr){4-5}
& $\eps$ & $\lvert\rho\rvert$ & $\eps$ & $\lvert\rho\rvert$ \\
\midrule
$0.25$ & $3.19$   & $0.9998$ & $0.263$ & $0.9998$ \\
$0.5$  & $6.39$   & $0.9998$ & $0.526$ & $0.9998$ \\
$1$    & $12.77$  & $0.9998$ & $1.052$ & $0.9998$ \\
$2$    & $25.55$  & $0.9998$ & $2.105$ & $0.9950$ \\
$4$    & $51.09$  & $0.9997$ & $4.209$ & $0.8785$ \\
$8$    & $102.19$ & $0.9997$ & $8.418$ & $0.7197$ \\
$16$   & $204.37$ & $0.9997$ & $16.84$ & $0.6848$ \\
$32$   & $408.74$ & $0.9997$ & $33.67$ & $0.6541$ \\
\bottomrule
\end{tabular}
\end{table}

\begin{figure}[t]
\centering
\includegraphics[width=\textwidth]{09_bandwidth_study.png}
\caption{Bandwidth sweep. On the RBF observations the diffusion map decays
monotonically from $0.9998$ onto the PCA baseline (dotted) as $\eps$ grows.}
\label{fig:bandwidth}
\end{figure}

The polynomial model is insensitive: $\lvert\rho\rvert$ moves by $10^{-4}$ across
a $128\times$ sweep, because a nearly straight manifold is ordered correctly at
any scale. The RBF model tells the real story: recovery is flat at $0.9998$ up to
$\times2$, then falls monotonically to $0.6541$ at $\times32$ ---
\textbf{converging onto PCA's $0.6624$}.

This is the mechanism behind the whole bandwidth question. As $\eps\to\infty$,
$K_{ij} = \exp(-D^2_{ij}/\eps) \approx 1 - D^2_{ij}/\eps$, so the kernel becomes
an affine function of the squared distances and the diffusion map degenerates to a
linear method operating on the Gram matrix --- that is, to PCA. \emph{An
overshooting bandwidth does not merely blur the embedding; it converts the
nonlinear method into the linear one it was supposed to beat.} Since the median
heuristic overshoots by $29\times$ on this model, it lands squarely in the
degenerate regime, and this is precisely why the preliminary configuration found
diffusion maps and PCA indistinguishable.

\subsection{Sensitivity to measurement noise}
\label{sec:res-noise}

\begin{table}[t]
\centering
\caption{Measurement-noise study at fixed bandwidth (Section~\ref{sec:bandwidth}).}
\label{tab:noise}
\begin{tabular}{@{}lrrrr@{}}
\toprule
& \multicolumn{2}{c}{polynomial} & \multicolumn{2}{c}{RBF} \\
\cmidrule(lr){2-3}\cmidrule(lr){4-5}
$\eta$ & DM $\lvert\rho\rvert$ & PCA $\lvert\rho\rvert$
       & DM $\lvert\rho\rvert$ & PCA $\lvert\rho\rvert$ \\
\midrule
$0.00$ & $1.0000$ & $1.0000$ & $1.0000$ & $0.6297$ \\
$0.05$ & $0.9999$ & $0.9999$ & $0.9999$ & $0.6450$ \\
$0.10$ & $0.9998$ & $0.9997$ & $0.9998$ & $0.6624$ \\
$0.20$ & $0.9992$ & $0.9987$ & $0.9987$ & $0.6847$ \\
$0.40$ & $0.9974$ & $0.9960$ & $0.9235$ & $0.7022$ \\
\bottomrule
\end{tabular}
\end{table}

\begin{figure}[t]
\centering
\includegraphics[width=\textwidth]{10_noise_study.png}
\caption{Robustness to measurement noise. The diffusion map's advantage on the
RBF observations persists across the whole range.}
\label{fig:noise}
\end{figure}

At $\eta=0$ both methods attain $\lvert\rho\rvert = 1.0000$ on the polynomial
model, consistent with Section~\ref{sec:identifiability}: the noiseless
observations are a deterministic injective function of $X$ and perfect monotone
recovery is attainable in principle. Degradation is graceful thereafter. The
diffusion map holds its RBF advantage across the entire sweep --- $0.9235$ against
$0.7022$ even at $\eta = 0.4$, where the noise is four times the default.

Before column standardization, measurement noise adds $2d\eta^2$ in expectation
to each squared pairwise distance. Standardization rescales that contribution, so
the expression is only a qualitative guide to why small isotropic noise is mild;
degradation becomes pronounced once noise competes with the intrinsic curve's
local separation. PCA's RBF
score \emph{rises} slightly with noise ($0.630\to0.702$), which is not an
improvement in any meaningful sense but an artifact of noise flattening the
folded hump of Figure~\ref{fig:recovery}(d) relative to its linear component.

\subsection{Sensitivity to barrier depth}
\label{sec:res-barrier}

\begin{table}[t]
\centering
\caption{Barrier study. The horizon is scaled to hold $T/\tau_{\mathrm{esc}}$
roughly fixed, giving $45$--$61$ committed transitions throughout. $1/k$ is the
Kramers escape time.}
\label{tab:barrier}
\begin{tabular}{@{}lrrrrrrrr@{}}
\toprule
& & & & \multicolumn{2}{c}{poly.\ well alignment}
& \multicolumn{2}{c}{RBF $\mu_2/\mu_1$} & \\
\cmidrule(lr){5-6}\cmidrule(lr){7-8}
$\sigma$ & $\Delta V/D$ & $T$ & trans.
& $\alpha=\tfrac12$ & $\alpha=1$
& $\alpha=\tfrac12$ & $\alpha=1$ & $1/k$ \\
\midrule
$0.45$ & $2.47$ & $2549$ & $50$ & $0.961$ & $0.900$ & $32.95$ & $4.75$ & $52.5$ \\
$0.50$ & $2.00$ & $1589$ & $45$ & $0.950$ & $0.838$ & $23.30$ & $5.11$ & $32.8$ \\
$0.55$ & $1.65$ & $1140$ & $49$ & $0.910$ & $0.834$ & $14.87$ & $4.66$ & $23.2$ \\
$0.70$ & $1.02$ & $600$  & $61$ & $0.841$ & $0.667$ & $10.17$ & $4.19$ & $12.3$ \\
\bottomrule
\end{tabular}
\end{table}

\begin{figure}[t]
\centering
\includegraphics[width=\textwidth]{11_barrier_study.png}
\caption{Barrier depth. The $\alpha=1/2$ and $\alpha=1$ curves converge as the
barrier deepens: when the point cloud is geometrically clustered, geometry and
dynamics agree.}
\label{fig:barrier}
\end{figure}

Two readings, and the second is the cleanest single result in this paper.

First, \textbf{the $\alpha$ effect is largest in the shallow-barrier regime}. On
the polynomial model the $\alpha=1/2$ to $\alpha=1$ gap in well alignment is
$0.061$ at $\Delta V/D = 2.47$ but $0.174$ at $\Delta V/D = 1.02$. When the
barrier is deep, points pile up near $\pm1$ and thin out between, so the sampled
set becomes geometrically two-clustered and even a density-free operator finds the
split. When the barrier is shallow, the barrier region is populated, the geometry
is a continuum, and only the density carries the two-state information --- which
$\alpha=1$ discards. Deep barriers hide the distinction; shallow ones expose it.
This retrospectively justifies keeping $\sigma=0.7$ as the primary setting.

Second, \textbf{the $\alpha=1$ spectral gap is insensitive to barrier depth}. Across a
$4.3\times$ range of Kramers times ($12.3$ to $52.5$) the RBF $\alpha=1$ ratio
sits at $4.75,\,5.11,\,4.66,\,4.19$ --- pinned near the geometric value $4$ and
showing no clear trend. Over the same range the $\alpha=1/2$ ratio moves
$10.17\to14.87\to23.30\to32.95$, tracking the trend in $1/k$ and lying within a
factor of $0.63$--$0.83$ of it. This is evidence that density-sensitive
normalization retains information about changing metastability. It is not a
direct recovery of the physical timescale: Kramers' formula is asymptotic here,
and the non-isometric observation map changes the operator's metric. The contrast
with the nearly fixed $\alpha=1$ geometric ratio is the robust conclusion.

\subsection{A metric that certifies its own failure}
\label{sec:degeneracy}

One negative result deserves its own statement, because it would silently
invalidate any study that used only the original three metrics.

An unconstrained kernel-sum scan can select a bandwidth at which the kernel is
numerically nearly reducible. The conservative MST scale for the polynomial
observations is $\eps_{\mathrm{MST}}=3.418$; at
$\eps=0.05\,\eps_{\mathrm{MST}}=0.171$ the leading eigenvalues round to
$(1,\,1,\,1,\,0.99992)$. The Gaussian graph remains connected at an affinity
cutoff of $10^{-12}$, but two extreme samples, $x\in[1.82,1.87]$, are so weakly
coupled that $\psi_1$ behaves like their component indicator: they sit near
$-0.70$, while the other $2998$ values agree to within $3.8\times10^{-8}$. The
coordinate carries no usable global information, and $2$-means isolates those two
points, giving balanced accuracy $0.5007$, essentially chance.

Its Spearman correlation with the true state is $\mathbf{-0.9997}$.

The mechanism is worth stating precisely, because it is not simply tied ranks.
Rank correlation is scale-free: it discards magnitude entirely and sees only
ordering. Inside the near-constant blob the residual variation --- some
$8$ orders of magnitude below the coordinate's full span of $0.71$, and
attributable to the smooth variation of the graph degree $d_i^{-1/2}$ along the
curve --- happens to be monotone in $x$. Spearman therefore locks onto numerical
dust and reports near-perfect recovery. \emph{Precisely the property that makes
Spearman the right headline metric for Q1 --- invariance to monotone
reparameterization, hence to scale --- is what makes it incapable of noticing
that the structure it is ranking is numerically negligible.}

The two metrics thus flatly contradict each other, and only one is useful. We
therefore report the number of computed eigenvalues within $10^{-8}$ of $1$ and
the fraction of distinct coordinate values. The implementation warns about a
disconnected or nearly reducible kernel without equating near-unit modes to a
literal component count. \textbf{Rank correlation must never be read from a
diffusion map without also inspecting the spectrum and coordinate scale.}

% =============================================================================
\section{Discussion}
\label{sec:discussion}
% =============================================================================

\subsection{Answers to the two questions}

\textbf{Q1 is answered affirmatively, with a condition.} The diffusion map
recovers the hidden state to $\lvert\rho\rvert = 0.9998$ through a nonlinear,
noisy, $15$-dimensional encoding, which by Section~\ref{sec:identifiability} is
the strongest form of recovery the problem admits. It beats PCA decisively
($0.9998$ against $0.6624$) when the observation manifold is strongly curved and
ties when it is not. The condition is the bandwidth: the advantage exists only
below roughly twice the scan-selected $\eps$ and vanishes smoothly onto the PCA
baseline above it.

\textbf{Q2 depends on both normalization and observation geometry.} On the
polynomial model, $\alpha=1$ gives well accuracy $0.5804$ while $\alpha\le1/2$
gives above $0.99$; on the RBF model, geometric well separation makes every
normalization score highly. Thus $\alpha$ controls how much sampling density is
retained, but the data geometry controls the practical size of the effect.

\subsection{Why the standard default is the wrong one here}
\label{sec:disc-geometry}

Two facts compound. The observation map is \emph{memoryless}: each $y_i$ depends
only on $X_i$, so shuffling the time index leaves the diffusion map exactly
unchanged, and the point cloud carries no information about ordering, correlation
times, or rates. The single imprint the dynamics leave is the sampling density,
since the particle dwells near the wells and points accumulate there. And
$\alpha=1$ is defined by the property that it divides that density out.

The data contain no dynamical information except the density; the normalization
deletes the density. Nothing goes wrong in the implementation --- the method
faithfully computes the operator it was asked to compute, and
Section~\ref{sec:res-spectra} confirms by the $k^2$ law that the operator is the
interval Laplacian. The mismatch is between that operator and the question.

This is consistent with the literature, with one important qualification. When
\citet{nadler2006diffusion} and \citet{coifman2008diffusion} connect diffusion
coordinates to reaction coordinates, the relevant normalization is $\alpha=1/2$.
It gives the backward Fokker--Planck generator only in coordinates with the
appropriate metric; under our nonlinear observation maps it is a
density-sensitive proxy rather than the exact generator of Eq.~\eqref{eq:sde}.
Choosing $\alpha=1$ --- the right default when the goal is to recover shape
independent of how densely one sampled --- silently changes the question from
``what are the slow variables of this process?'' to ``what does this set of
points look like?''. For Q1 the latter suffices. For Q2 it discards the one
dynamical imprint available to a memoryless point cloud, although geometry may
still separate the wells as it does for the RBF encoding.

\subsection{Why the preliminary configuration was misleading}
\label{sec:disc-pilot}

The earlier configuration reported $\lvert\rho\rvert = 0.9992$ and $0.9846$ well
accuracy and concluded that diffusion maps recover the state, tie with PCA, and
separate the wells. The first claim was right, the second was an artifact, and
the third was false. Four defects (Table~\ref{tab:corrections}) each contributed,
and none was individually unreasonable:

\begin{itemize}
\item \textbf{The horizon was too short.} $T=30$ is $2.4$ escape times; the path
made $2$--$3$ committed transitions and spent $84.9\%$ of its time in one well.
No dynamical quantity was estimable, and the imbalance meant a $2$-means split of
\emph{any} monotone coordinate landed near $x=0$, manufacturing the $0.9846$. With
balanced occupancy the same configuration yields $0.5804$ --- the reported number
moved by $40$ points once the sampling was fixed, with no change to the method.

\item \textbf{The bandwidth overshot.} The median heuristic sat $29\times$ above
the scan-selected value on curved data, i.e.\ deep in the regime where the
diffusion map degenerates to PCA (Section~\ref{sec:res-bandwidth}). The observed
tie with PCA was therefore not a finding about manifold learning but a
consequence of the hyperparameter.

\item \textbf{Only $\alpha=1$ was run.} This density-free normalization suppresses
the well structure on the compact polynomial arc, so the decisive normalization
comparison was never made.

\item \textbf{Only one observation model was used}, and its image was nearly
straight, so PCA had no reason to fail and the comparison carried no information.
\end{itemize}

The general lesson is that the defects were mutually camouflaging. The short
horizon inflated the well score, which made the $\alpha=1$ failure invisible; the
overshooting bandwidth produced a PCA tie, which looked like a substantive
negative result rather than a hyperparameter artifact; and the single
near-linear observation model gave no contrast that would have exposed either.
Each check that would have caught one defect was itself compromised by another.

\subsection{Practical guidance}

\begin{enumerate}
\item \textbf{Choose $\alpha$ from the question.} Use $\alpha=1$ for geometry
independent of sampling density and $\alpha=1/2$ when density-sensitive structure
is relevant. Recovering the exact physical generator additionally requires the
correct metric or a dynamics-tailored kernel.

\item \textbf{Verify which operator you converged to.} Gap ratios near
$(4,9,16)$ support the interval Laplacian --- geometry. A large
$\mu_2/\mu_1$ that moves with a physical parameter is evidence of retained
density information, not by itself proof that the exact physical generator was
recovered.

\item \textbf{Check conditioning before reading any score.} Rank correlation can
report $0.9997$ for a numerically nearly reducible kernel
(Section~\ref{sec:degeneracy}).

\item \textbf{Do not trust a bandwidth heuristic based on typical distances.}
Overshooting converts the diffusion map into PCA. Sweep $\eps$ and look for the
plateau; if the method's advantage over a linear baseline disappears at large
$\eps$, that is the expected behaviour, not evidence about the method.

\item \textbf{Test manifold learning against PCA on a curved manifold.} A pipeline
can look aggressively nonlinear and still present a nearly linear problem, since
what matters is the linearity of the \emph{leading} direction, not the
nonlinearity of individual coordinates.
\end{enumerate}

% =============================================================================
\section{Limitations and Future Work}
\label{sec:limitations}
% =============================================================================

\paragraph{Limitations.}
(i) Every result rests on a single trajectory from a single seed; no error bars
over simulation replicates are reported, so the small differences --- the
polynomial-model DM/PCA margins, PCA's non-monotone RBF noise trend --- should
not be interpreted. The large effects ($0.9998$ versus $0.6624$; $0.9960$ versus
$0.5804$) are unlikely to reverse, but their seed variability has not been
measured.
(ii) The automatic bandwidth selector is not reliable
(Section~\ref{sec:bandwidth}); the bandwidth sweep, not the selector, carries the
evidence.
(iii) Both observation models are memoryless, so no method operating on this
point cloud can access dynamical information beyond the invariant density.
(iv) The dense $O(N^2)$ formulation caps $N$ at a few thousand, which is why the
long horizon had to be bought by subsampling rather than by more points.
(v) The barrier study varies $\sigma$, which changes both the barrier depth and
the width of the sampled region; the two effects are not separated.
(vi) The two-cluster variance score is descriptive rather than a formal modality
test, and the well-alignment score uses ground truth; neither alone proves
dynamical recovery.
(vii) The Kramers comparisons in Sections~\ref{sec:res-trajectory}
and~\ref{sec:res-barrier} use an asymptotic formula at $\Delta V/D$ between $1$
and $2.5$, and agree in trend and order of magnitude but not in constant.

\paragraph{Future work.}
The clearest next step is \textbf{time-lagged kernels}. Every limitation of type
(iii) follows from the observations being memoryless; a kernel built on
delay-embedded observations, or TICA \citep{perez2013identification}, uses the
time-lagged covariance and therefore sees the dynamics directly rather than
through the density. That would let the reaction coordinate be recovered without
relying on $\alpha$ to reintroduce information the encoding discarded. Beyond it:

\begin{itemize}
\item \textbf{It\^o-tailored kernels.} The construction of
\citet{banisch2017diffusion} targets the generator of an arbitrary
non-degenerate It\^o process directly, and would apply unchanged to irreversible
or state-dependent-noise variants of Eq.~\eqref{eq:sde}. Even in the present
problem, standard $\alpha=1/2$ normalization does not by itself correct the metric
distortion introduced by a nonlinear observation map.
\item \textbf{Variable bandwidth.} Locally scaled
\citep{rohrdanz2011determination} or variable-bandwidth
\citep{berry2016variable} kernels, which would address both the selector's
instability and the non-uniform sampling that causes it.
\item \textbf{Sparse and Nystr\"om methods} \citep{fowlkes2004spectral} to lift
the $O(N^2)$ cap and allow long horizons without subsampling.
\item \textbf{Replicate seeds} to put error bars on every table.
\item \textbf{Higher-dimensional and asymmetric potentials}, where the reaction
coordinate is a genuinely nonlinear function of the observables and the
transition path is curved.
\end{itemize}

% =============================================================================
\section{Conclusion}
\label{sec:conclusion}
% =============================================================================

We implemented diffusion maps from first principles and asked what they recover
from a stochastic double-well system observed through a nonlinear, noisy,
$15$-dimensional encoding. The answer separates into two.

On \emph{state recovery}, diffusion maps work and beat the linear baseline
exactly where theory says they should. The leading nontrivial coordinate orders
the samples as the hidden state does to
$\lvert\rho_{\mathrm{Spearman}}\rvert = 0.9998$ on both observation models, while
PCA manages only $0.6624$ on the curved one, where its first component folds into
a non-monotone hump. The advantage is real but conditional: as the bandwidth
grows the Gaussian kernel becomes affine in the squared distances and the
diffusion map degenerates onto PCA, which we observe directly as recovery decays
from $0.9998$ to $0.654$ against PCA's $0.662$.

On \emph{metastable structure}, normalization and observation geometry interact.
Because the observations are memoryless, sampling density is the only dynamical
imprint available to the point cloud, and $\alpha=1$ removes it. On the polynomial
model, well accuracy consequently falls from $0.9960$ to $0.5804$ as $\alpha$
increases; on the geometrically separated RBF model it remains high. At
$\alpha=1$ the RBF gap ratios sit near the interval $k^2$ law and remain nearly
fixed as the barrier changes, while at $\alpha=1/2$ the leading gap grows with the
Kramers-time trend. The latter is evidence of retained density information, not
an exact recovery of the physical generator, because the nonlinear observation
map changes the metric.

We also documented how a plausible-looking configuration of this same study
reached the opposite conclusions, and why: a horizon of $2.4$ escape times made a
$2$-means threshold land near the barrier by luck and manufacture a $0.9846$ well
accuracy that is $0.5804$ once the sampling is balanced; a median-distance
bandwidth overshot by $29\times$ into the regime where the method is PCA; and
rank correlation reported $0.9997$ for a nearly reducible kernel whose leading
eigenvector behaved like a two-point component indicator. The defects were mutually
camouflaging, and none was individually unreasonable.

The transferable conclusion is narrow and practical. A diffusion map answers
whichever question its normalization encodes, at whichever scale its bandwidth
selects, and its headline metric can certify an embedding that carries no
information at all. Before interpreting slow modes as reaction coordinates, it is
worth inspecting near-unit eigenvalues, the $k^2$ signature, and a bandwidth
sweep, and remembering that the observation metric determines which continuum
operator is approximated.

% =============================================================================
\section*{Reproducibility Statement}
% =============================================================================
\addcontentsline{toc}{section}{Reproducibility Statement}

Every reported number and table is derived from the CSV outputs of
\texttt{python run\_experiment.py}; the numerical values are transcribed into the
\LaTeX\ source, while the figures are generated directly by the script. The run
uses the parameters of Table~\ref{tab:params}, with no network access or external
data. Outputs are written to \texttt{outputs/data/}, \texttt{outputs/metrics/}
(four per-study CSVs, a combined table, and the resolved parameters as JSON), and
\texttt{outputs/figures/} (seventeen 300-DPI PNGs). The degeneracy demonstration
of Section~\ref{sec:degeneracy} is reproduced by the script in
Appendix~\ref{app:degeneracy}. Outputs of the superseded preliminary
configuration are retained under \texttt{outputs/\_superseded/} for comparison.
The test suite is run with \texttt{pytest} ($55$ passed).

% =============================================================================
\bibliographystyle{plainnat}
\bibliography{references}

\appendix
\newpage

% =============================================================================
\section{Notation}
\label{app:notation}
% =============================================================================

\begin{center}
\begin{tabular}{@{}ll@{}}
\toprule
Symbol & Meaning \\
\midrule
$X_t,\ X_n$ & hidden scalar state; continuous and discrete time \\
$V(x)$ & double-well potential $-x^2/2 + x^4/4$ \\
$\sigma$, $D$ & process-noise coefficient; diffusion constant $\sigma^2/2$ \\
$\Delta V$, $k$ & barrier height $1/4$; Kramers escape rate \\
$\Delta t$, $N$, $T$ & integration step, retained samples, horizon \\
$\rho_s$ & stationary density $\propto e^{-2V/\sigma^2}$ \\
$\phi(x)$, $m$ & feature map and its dimension ($9$ polynomial, $12$ RBF) \\
$A$, $d$ & random lift matrix; ambient dimension ($15$) \\
$\eta$ & measurement-noise standard deviation \\
$Y$, $y_i$ & standardized observation matrix and its rows \\
$\mathcal{M}$ & the intrinsic one-dimensional observation curve \\
$\eps$ & kernel bandwidth (a squared length scale) \\
$K$, $\Kt$ & Gaussian kernel; $\alpha$-normalized kernel \\
$q_i$, $d_i$, $D$ & density proxy, graph degree, degree matrix \\
$\alpha$ & anisotropic normalization exponent \\
$P$, $P_{\mathrm{sym}}$ & Markov matrix and its symmetric conjugate \\
$\lambda_k$, $\psi_k$ & eigenvalues and right eigenvectors of $P$ \\
$\mu_k$ & continuum eigenvalues; $\mu_k/\mu_1 \approx
          (1-\lambda_k)/(1-\lambda_1)$ \\
$\mathrm{PC}_1$ & first principal component of $Y$ \\
\bottomrule
\end{tabular}
\end{center}

% =============================================================================
\section{Reproducing the Degeneracy Demonstration}
\label{app:degeneracy}
% =============================================================================

Run from the project root after \texttt{python run\_experiment.py}. It reproduces
the numbers of Section~\ref{sec:degeneracy}: a nearly reducible kernel whose
leading eigenvector behaves like a two-point component indicator, yet whose
Spearman correlation with the true state reads $-0.9997$.

\begin{small}
\begin{verbatim}
import warnings
import numpy as np, sys; sys.path.insert(0, '.')
from scipy.stats import spearmanr
from src.diffusion_maps import (DiffusionMapConfig, compute_diffusion_map,
                                connectivity_scale, squared_distance_matrix)
from src.evaluation import well_classification_score

Y = np.load('outputs/data/observations_poly.npy')
x = np.load('outputs/data/latent_state.npy')

mst_scale = connectivity_scale(squared_distance_matrix(Y))
print(f'MST scale         = {mst_scale:.4f}')

with warnings.catch_warnings():
    warnings.simplefilter('ignore')                  # we are provoking it
    r = compute_diffusion_map(Y, DiffusionMapConfig(epsilon=0.05 * mst_scale,
                                                    n_components=3))
p1 = r.coordinates[:, 0]
print(f'eps              = {r.epsilon:.4f}')
print(f'eigenvalues      = {np.round(r.eigenvalues[:4], 6)}')
print(f'near-unit modes  = {r.n_trivial}   degenerate = {r.is_degenerate}')
print(f'psi_1 full range = {p1.max() - p1.min():.4f}')
print(f'psi_1 IQR        = {np.subtract(*np.percentile(p1, [75, 25])):.3e}'
      f'   <-- 8 orders below the range')
print(f'Spearman with X  = {spearmanr(p1, x).correlation:+.4f}   <-- numerical dust')
print(f'well accuracy    = {well_classification_score(p1, x):.4f}   <-- chance')
\end{verbatim}
\end{small}

% =============================================================================
\section{Output Manifest}
\label{app:manifest}
% =============================================================================

\begin{small}
\begin{verbatim}
outputs/
  data/
    time.npy, latent_state.npy         # retained time grid and state
    latent_state_full.npy              # unsubsampled path (transition stats)
    observations_{poly,rbf}.npy        # standardized Y for each model
    features_{poly,rbf}.npy            # feature matrices F
    dm_coordinates_{poly,rbf}.npy      # diffusion coordinates
    pca_scores_{poly,rbf}.npy          # principal-component scores
  metrics/
    alpha_study.csv                    # normalization study
    bandwidth_study.csv                # bandwidth study
    measurement_noise_study.csv        # measurement-noise study
    barrier_study.csv                  # barrier-depth study
    combined_metrics.csv               # all four, tagged by 'study'
    run_parameters.json                # resolved run parameters
  figures/
    01_state_vs_time.png  02_state_histogram.png
    03_{poly,rbf}_feature_traces.png
    04_{poly,rbf}_dm_coordinate_vs_state.png
    05_{poly,rbf}_pca_coordinate_vs_state.png
    06_{poly,rbf}_dm_embedding_scatter.png
    07_{poly,rbf}_pca_embedding_scatter.png
    08_alpha_study.png    09_bandwidth_study.png
    10_noise_study.png    11_barrier_study.png
    12_alpha_spectra.png
  _superseded/                         # outputs of the preliminary run
\end{verbatim}
\end{small}

\end{document}
